\paragraph{Letter to Dieudonné, 29 September 1965 (cote 25, sites 134--135).}
The handwritten date fixes this letter to 29 September 1965. Grothendieck thanks Dieudonné for a contents list for paragraphs 16--19 and asks for the provisional contents of 20 and 21, accepting their inclusion in the fourth fascicle of Chapter IV while asking how his former paragraph 20 will be divided. The proposed continuation assigns linear systems and additions on the Picard group to 22, Grassmannians to 23, smooth forms and ordinary quadratic singularities to 24, hyperplane sections to 25, the resultant and discriminant to 26, and infinitesimal extensions to 27; 20 and 21 are still represented by question marks. His joke about the pleasing number $27=3^3$ accompanies an expectation that the long paragraph on hyperplane sections will be divided again. These are the numbers of the plan recorded by this letter. The scan supplies no identification of them with every later number used in the surrounding prenotes. The second subject is authorship: he declines to publish the former appendix to paragraph 18 under his name alone, credits Dieudonné's extensive rewriting, and asks him to reconsider a joint paper. Site 135 also preserves a completely canceled paragraph referring to ``par.20, 10.9.1.'' and to paragraph 12: it invokes constructibility, and then openness under the stated flatness and finite-presentation assumptions, of a fiberwise depth locus, as part of the familiar argument for passing farther along the indexing system. The handwritten deletion belongs to the letter's textual history and remains visible in its diplomatic form. The check confirms the date, the cancellation, and the proposed contents directly from the two scans; it establishes no bibliographic identification with the separately named L49d. See \href{https://grothendieck.umontpellier.fr/web/viewer.html?file=../25.pdf#page=135}{cote 25, site 134} and \href{https://grothendieck.umontpellier.fr/web/viewer.html?file=../25.pdf#page=136}{site 135}.
